\SourceLeaf{380}%
eine Zeit lang gestört, und der Kollimationsfehler schwankte beträchtlich, wie es zuvor \textit{nie} vorgekommen war. Zunächst wurde die Ursache vergeblich in mancherlei anderen Umständen gesucht. Sie fand sich in einem kleinen Knoten des Fadens, der um die den Verschluss der Libelle sichernde Blasenhaut gebunden war. Der Knoten war ein klein wenig zu dick und berührte die Innenfläche der Hülse. Nachdem er weggeschnitten war, sodass die Glasröhre nur noch die leicht vorstehenden Schraubenspitzen berührte, waren die Beständigkeit des Kollimationsfehlers und die frühere schöne Übereinstimmung aller Beobachtungen sogleich wiederhergestellt. Auch bei den hier behandelten Turiner Beobachtungen bemerken wir beträchtliche Schwankungen des Kollimationsfehlers. Wir meinen die kleineren, die zufällig scheinen. Die größeren Änderungen um mehrere Minuten hatten ohne Zweifel ihren guten, dem Astronomen bekannten Grund, auch wenn er bei den Beobachtungen nicht genannt ist. Heute, da man ein so vorzügliches Mittel hat, den Kollimationsfehler jederzeit ohne Umlegen zu bestimmen, lässt sich die Ursache ähnlicher Anomalien umso leichter finden.

\begingroup\rightskip=0pt plus 1em
Der vierte Abschnitt beschreibt den Anschluss des Mont Cenis an das Dreiecksnetz durch eine besondere Triangulierung und eine kurze, auf dem Bergplateau gemessene Basislinie sowie die astronomische Bestimmung des Azimuts der Verbindungslinie Mont Cenis --- Bellecombe. Letztere wurde zweimal vorgenommen. Die Ergebnisse aus den Jahren 1821, 1822, mit Repetitionskreisen von \textsc{Throughton} und \textsc{Reichenbach}, weichen um \(8''6\) voneinander ab. Man nahm den Mittelwert aus beiden, obwohl die spätere Bestimmung bei Weitem zuverlässiger scheint.\par
\endgroup

Ebenso enthalten die beiden folgenden Abschnitte die astronomischen Bestimmungen der Azimute der Richtungslinien Mont Colombier --- Mont Granier und Turin neue Sternwarte --- Superga. In allen drei Fällen dienten die Meridianzeichen der jeweiligen Passagenfernrohre als Grundlage für diese Bestimmungen. Schließlich finden sich im sechsten Abschnitt noch die Arbeiten, mit denen eine frühere astronomische Azimutbestimmung von \textsc{Oriani} an der Mailänder Sternwarte auf die Orientierung der Seite Mailand Domturm --- Busto im französischen Dreiecksnetz übertragen wurde.

Im siebten Abschnitt werden nun aus diesen umfangreichen Arbeiten die Ergebnisse für die Längengradmessung abgeleitet. Man bezog die Messungen
